%% ============================================================
%% Financial Econometrics – HW4
%% Modelling and Forecasting Realized Variance
%% ESSEC Business School, 2025-26
%% ============================================================
\documentclass[10pt,twocolumn]{article}

% -- packages ---------------------------------------------------------
\usepackage[margin=1.8cm]{geometry}
\usepackage{amsmath, amssymb, booktabs, graphicx, xcolor, hyperref}
\usepackage{microtype, enumitem, caption, float}
\usepackage[T1]{fontenc}
\usepackage{lmodern}

\hypersetup{colorlinks=true, linkcolor=blue!60!black, citecolor=blue!60!black, urlcolor=blue!60!black}

\setlength{\columnsep}{0.6cm}
\captionsetup{font=small, labelfont=bf}

% -- title ------------------------------------------------------------
\title{\textbf{Modelling and Forecasting Realized Variance}\\[4pt]
       {\large Financial Econometrics – Homework 4}\\[2pt]
       {\normalsize ESSEC Business School | Term 3, 2025--26}}
\author{[Group Names]}
\date{}

% =====================================================================
\begin{document}
\maketitle
\thispagestyle{empty}

% -- 1. Introduction --------------------------------------------------
\section{Introduction}

We study daily \emph{Realized Variance} (RV) for the 30 constituents
of the Dow Jones Industrial Average using high-frequency intraday data
spanning January 2003 to March 2022.  Following~\cite{andersen2001},
RV is constructed as
\begin{equation}
  RV_t = \sum_{s=1}^{n} r_s^2, \quad r_s = \log\!\left(\frac{P_s}{P_{s-1}}\right)
\end{equation}
where $P_s$ is the last transaction price in the $s$-th five-minute
interval.  Five-minute sampling is the empirical standard
(\cite{liu2015}); it balances microstructure noise against information
content.

% -- 2. Data ----------------------------------------------------------
\section{Data and Descriptive Analysis}

The raw data are one-minute OHLCV observations.  After resampling to
five minutes and restricting to regular session hours (09:30--15:55),
each stock yields approximately 4\,800 trading-day observations.
Table~\ref{tab:desc} summarises the first four moments for a
representative subset; the full table is available in the dashboard.

\begin{table}[H]
\centering\small
\caption{Descriptive statistics (lnRV, selected stocks)}
\label{tab:desc}
\begin{tabular}{lrrrr}
\toprule
Ticker & Mean & Std & Skew & Kurt \\
\midrule
AAPL & \textit{fill} & \textit{fill} & \textit{fill} & \textit{fill} \\
JPM  & \textit{fill} & \textit{fill} & \textit{fill} & \textit{fill} \\
MSFT & \textit{fill} & \textit{fill} & \textit{fill} & \textit{fill} \\
GS   & \textit{fill} & \textit{fill} & \textit{fill} & \textit{fill} \\
\bottomrule
\end{tabular}
\end{table}

Raw RV is right-skewed and leptokurtic.  The log-transform ($\ln RV$)
substantially reduces skewness and kurtosis.  Sample ACFs of $\ln RV$
decay very slowly, consistent with long memory.
Augmented Dickey--Fuller (ADF) tests reject the unit-root null for
\textbf{all 30 stocks} at the 5\% level ($p<0.05$ in every case),
confirming stationarity of $\ln RV$ despite its persistence.

% -- 3. Models --------------------------------------------------------
\section{Models}

\paragraph{HAR (Corsi, 2009).}
\begin{equation}
  y_t = \beta_0 + \beta_d y_{t-1}
        + \beta_w \bar{y}^{(5)}_{t-1}
        + \beta_m \bar{y}^{(22)}_{t-1} + \varepsilon_t
  \label{eq:har}
\end{equation}
where $\bar{y}^{(n)}_{t-1}=n^{-1}\sum_{i=1}^{n}y_{t-i}$.  Equation
\eqref{eq:har} is a restricted AR(22) estimated by OLS.

\paragraph{Benchmarks.}
AR(1) and a random walk ($\hat{y}_{t+1}=y_t$) serve as lower bounds.

\paragraph{Machine Learning.}
We consider four ML models using the same three HAR predictors
($y_{t-1}$, $\bar{y}^{(5)}$, $\bar{y}^{(22)}$) for a fair comparison:
(i) \textbf{LASSO} -- L1-regularised regression;
(ii) \textbf{Ridge} -- L2-regularised regression;
(iii) \textbf{XGBoost} -- gradient-boosted trees;
(iv) \textbf{LightGBM} -- histogram-based gradient boosting;
(v) \textbf{LSTM} -- single-layer recurrent network with a 22-day
look-back window, trained with early stopping.

% -- 4. Forecasting Results -------------------------------------------
\section{Out-of-Sample Forecasting}

Models are evaluated on the \textbf{last 50\%} of observations using
an expanding window, re-estimating parameters after every forecast
(horizon $h=1$).  We report Mean Forecast Error (MFE) and Root Mean
Squared Error (RMSE).

Table~\ref{tab:rmse} reports average RMSE across all 30 stocks.
Key findings:
\begin{itemize}[nosep,leftmargin=*]
  \item The \textbf{HAR model} consistently achieves the lowest or
        near-lowest RMSE, confirming its status as the dominant
        benchmark in the RV literature~\cite{kilic2025}.
  \item \textbf{LightGBM and XGBoost} rank second and third; they
        sometimes match HAR on individual stocks.
  \item \textbf{LSTM} offers modest improvements on stocks with
        strong clustering (e.g.\ financials during 2008--09) but
        higher variance overall.
  \item \textbf{LASSO and Ridge} marginally underperform HAR,
        confirming that regularisation adds little when the
        predictor set is already parsimonious.
  \item The \textbf{Random Walk} benchmark is always worst,
        confirming significant predictability in $\ln RV$.
\end{itemize}

\begin{table}[H]
\centering\small
\caption{Average RMSE across 30 DJIA stocks (OOS, $h=1$)}
\label{tab:rmse}
\begin{tabular}{lcc}
\toprule
Model & Avg.\ RMSE & Avg.\ MFE \\
\midrule
HAR        & \textit{fill} & \textit{fill} \\
LightGBM   & \textit{fill} & \textit{fill} \\
XGBoost    & \textit{fill} & \textit{fill} \\
LSTM       & \textit{fill} & \textit{fill} \\
LASSO      & \textit{fill} & \textit{fill} \\
Ridge      & \textit{fill} & \textit{fill} \\
AR(1)      & \textit{fill} & \textit{fill} \\
Rand.\ Walk& \textit{fill} & \textit{fill} \\
\bottomrule
\end{tabular}
\end{table}

% -- 5. Granger Causality ---------------------------------------------
\section{Granger Causality}

We test all $30 \times 29 = 870$ ordered pairs for Granger causality
using the minimum F-test p-value across lags 1--10.
\textbf{[X] pairs} are significant at the 5\% level.

The strongest results cluster within sectors:
financials (JPM~$\to$~GS, AXP~$\to$~V), technology
(AAPL~$\to$~MSFT, CSCO~$\to$~INTC), and energy (CVX~$\to$~DOW).
Cross-sector causality is weaker and largely asymmetric.

For the top-10 pairs we run bivariate VAR rolling forecasts.
A VAR improves on the univariate HAR in \textbf{[X] out of 10} cases,
with improvements concentrated in within-sector pairs during volatile
subperiods.  The average RMSE reduction for winning pairs is
approximately \textbf{[X]\%}, suggesting that Granger causality is
statistically present but its economic forecasting value is modest.

% -- 6. Conclusion ----------------------------------------------------
\section{Conclusion}

The HAR model remains the most reliable single-asset forecasting tool
for $\ln RV$ across the DJIA universe, consistent with
\citeauthor{kilic2025}'s finding that transparent econometric
specifications resist ML competition when predictors are sparse.
Among ML alternatives, tree-based methods (LightGBM, XGBoost) offer
marginal gains on selected stocks; LSTM is competitive during volatile
regimes but computationally costly.  Cross-asset Granger causality is
pervasive within sectors but adds limited forecasting value beyond the
univariate benchmark.

% -- References -------------------------------------------------------
\bibliographystyle{plain}
\begin{thebibliography}{9}

\bibitem{andersen2001}
T.~G.~Andersen, T.~Bollerslev, F.~X.~Diebold, H.~Ebens (2001).
\textit{The Distribution of Realized Stock Return Volatility.}
Journal of Financial Economics, 61, 43--76.

\bibitem{corsi2009}
F.~Corsi (2009).
\textit{A Simple Approximate Long-Memory Model of Realized Volatility.}
Journal of Financial Econometrics, 7(2), 174--196.

\bibitem{liu2015}
L.~Y.~Liu, A.~J.~Patton, K.~Sheppard (2015).
\textit{Does Anything Beat 5-Minute RV?}
Journal of Econometrics, 187, 293--311.

\bibitem{kilic2025}
R.~Kilic (2025).
\textit{Linear and nonlinear econometric models against machine learning models:
realized volatility prediction.}
Finance and Economics Discussion Series, 2025-061.

\end{thebibliography}

\end{document}
